% Transcription — fonds Alexandre Grothendieck, shelfmark 148, batch 4 (pages 61-71).
% Established from the facsimile archives/batches/148/batch-04.pdf, cut from a
% copy of 148.pdf fetched through the project's relay
% (https://grothendieck-relay.onrender.com/source/148.pdf); PDF page k is the
% archivists' page 60+k, no cover sheet. Per the skill transcribe-grothendieck.
% The folder is seventy-one pages in four batches; this is the last (eleven
% pages). Batches 1-3 are transcribed in parallel by other agents; page 60
% (his p. 17, « Récapitulation des opérations élémentaires » 1°-7°) was looked
% at, at low resolution, only to fix the numbering of the operations.
%
% Pass: Opus 5.5 (claude-opus-5-5), 2026-09-26 — first pass, unchecked against the pages by a human.
%
% Pages skipped: 65 only, as a page: it is a cover (tan listing paper, blank
% but for three pencil lines of his at top right, « Découpage de cartes en
% "pantalons" généralisés »); by the rule for covers it takes no \page and his
% words become the section heading, with a \note. Pages summarised: none (no
% typescript). Pages 67, 68, 69 and 71 are mostly or wholly drawings; they are
% kept, their drawings described in \note{}s and every word or formula of his
% transcribed. Page 67 is the back of a computer listing (upside down, printed
% « 02 JUN 82 ») carrying only his pencil sketches: kept because the sketches
% are his and belong to the pantalon run; the printed date is recorded once in
% a note, as the skill asks, and is not used in \dating.
%
% His own pagination: pages 61-64 carry his numbers 18, 19, 20, 21 (top
% centre), continuing a run from batch 3 (p. 60 = his 17); the run ends
% mid-page 64 with « ... ». Pages 65-71 carry no number of his. He refers to
% the elementary operations of his p. 17 by their numbers 1°-7° (4° and 5°
% circled on pp. 61, 63; kept as (4°), (5°) with a note). His formula mark
% (**) on p. 62; (*) on p. 61 in a cancelled diagram. On p. 70 his own
% numbering 1), 2°), 3°) and « Variantes 1°) ». No dates in his hand anywhere
% in the batch.
%
% What the batch is about. Pages 61-64 (his 18-21): a programme for a
% « petit théorème de dévissage » — the topological Teichmüller groupoid
% S_f T D^±_top, with its functors to weighted graphs Γ and to [±1] and the
% elementary operations 1°-7° (except 4° and 5°), should be reconstructible
% from the groupoids S_1 T^±_{g,ν} -> T^±_{g,ν} -> Ens(ν) × [±1]; S_1 T^±_{g,ν}
% defined as the (X, R) with X connected of genus g, ν boundary components,
% one point of R on each, i.e. weighted graph a star with ν edges of weight 1.
% The essential question: reconstruct T^±_top(Γ) functorially in Γ and make
% the operations explicit; the operation 4° will come with the description,
% the « bouclages de trous » via the special cases for ST_{g,ν}, T_{g,ν}, to
% be put among the « elementary » structures, e.g. the diminution (**)
% T^±_{g,ν,ν'} -> T_{g,ν?ν'}. But bouclages are superficial: to relate genus g
% to g' < g one needs gommage of cutting circles (5° as he writes it here);
% further, the fundamental groupoids Π_1(X; R) of the (X, [Σ~], R) should be
% made explicit, functorially in R, and a new object of groupoids over ST^±_top
% — already non-trivial for S_2 T_{0,3}, perhaps too ambitious (a « modèle »
% containing all the cords « ça fait un peu gros ! »). Page 65: cover
% « Découpage de cartes en "pantalons" généralisés ». Page 66 (a sketch sheet,
% partly pencil, partly ink): a standard subdivided pantalon with its three
% holes <-> a « triangle pondéré » J -> N; a map cut « pantalonniquement »
% <-> a generalised pantalon map (trivalent, weighted half-edges, α_s+β_s =
% α_t+β_t, anti-isomorphisms of torsors under Z/γ, circular sums); the
% relations β+γ = α', γ+α = β', α+β = γ' solved for α, β, γ (NB integrality iff
% α'+β'+γ' even); a count for the sphere with three vertices: α+β+γ edges,
% (α−1)+(β−1)+(γ−1) bigons and 2 triangles, α+β+γ−1 faces; a drawing of the
% three-holed sphere (holes 0, 1, ∞) subdivided, with points a_i, b_i labelled
% by hole. Pages 67-69, 71: drawings of maps on the sphere with nested
% « pantalon » curves (red ink over pencil); p. 69 carries a square of
% inclusions T_top ⊂ TS_top, T^±_top ⊂ STS^±_top (« imm ») and « fid. ess.
% surj. » verticals, π_0(T_top) ≃ π_0(T^±_top) ≃; p. 71 is captioned
% « découpage pantalonnesque d'une carte (triangulée) » with a weighted
% trivalent graph. Page 70: « Groupoïdes de Teichmüller » — definitions: 1) ST_top,
% objects (X, S), X an oriented compact surface with boundary, S a finite part
% of the boundary, arrows isotopy classes of isomorphisms (« groupoïde de
% Teichmüller spécial orienté topologique »); S = ∅ gives T_top; variant
% ST^±_top (orientation not respected), with a functor to [±1] whose fibre is
% ST_top; 2°) (X, Σ, S) with Σ a compact 1-dimensional submanifold of X − ∂X,
% S ⊂ Σ ∪ ∂X finite: TD_top, D the initial of « divisage » (« groupoïde de
% Teichmüller orienté divisé topologique »); NB ST_top ⊂ STD_top open immersion,
% the subgroupoid Σ = ∅; 3°) likewise TS^±_top (« biorienté divisé »).
%
% Notation fixed once for the batch: his script T as \mathcal{T}; the prefix S
% (« spécial », p. 70) and the suffixes D (« divisé »), S, and the index
% top as \mathrm{top}; ± as a superscript; [±1] for the groupoid of the group
% ±1; Π_1 for the fundamental groupoid; Γ underlined (a weighted graph) as
% \underline{\Gamma}; Σ~ as \widetilde{\Sigma}. The small index under the S
% on p. 61 (read f, perhaps another letter) is flagged at its first occurrence.
%
% Hardest pages: 63 (the long sentence on Π_1 and the « nouvel objet », with a
% cramped left-margin note), 61 (a cancelled first diagram and a margin note
% written at 45°), 66 (a sketch sheet with the text block overlaid on
% drawings). Readings to check first: « La question essentielle ... » (p. 62),
% the sign between ν and ν' in (**) (p. 62), the margin notes of pp. 61, 62,
% 63, the words around « revêtements » on p. 63, the parenthesis on p. 64,
% the top block of p. 66, and the index of S on p. 61.
\documentclass[11pt,a4paper]{article}
% Le préambule commun à toutes les transcriptions du fonds Grothendieck.
%
% Il tient en une idée : **l'incertitude se compose, elle ne se résout pas.**
% Une transcription qui lisse un mot douteux a détruit la seule chose pour
% laquelle on la faisait — un lecteur doit pouvoir savoir, à chaque endroit,
% ce qui était sur le feuillet et ce qui a été ajouté. Les six macros qui
% suivent sont donc l'appareil critique, et non de la décoration.
%
% Les mêmes commandes sont reconnues par `scripts/render.mjs`, qui produit la
% vue de lecture du site. Une macro ajoutée ici doit l'être aussi là-bas,
% faute de quoi le rendu échouera bruyamment — ce qui est voulu : un rendu
% silencieusement faux est pire qu'un rendu absent.

\ProvidesPackage{grothendieck}[2026/08/08 Transcriptions du fonds Grothendieck]

\usepackage{amsmath,amssymb,amsthm}
\usepackage{mathrsfs}
\usepackage{stmaryrd}
% Les diagrammes commutatifs sont partout dans ce fonds : c'est la langue
% même de la géométrie algébrique de Grothendieck, et une transcription qui
% les rendrait en prose ne transcrirait rien.
\usepackage{tikz-cd}
% Français par défaut — c'est la langue des feuillets — mais l'anglais est
% chargé aussi : la traduction et le résumé sont des documents anglais, et la
% typographie française (espaces avant les deux-points) y serait une faute.
% Ces documents basculent par \selectlanguage{english} en tête de corps.
\usepackage[english,french]{babel}
\usepackage{fontspec}
\usepackage{geometry}
\geometry{a4paper,margin=25mm,marginparwidth=28mm}
\usepackage{marginnote}
\usepackage{xcolor}
% ulem, not soul. `soul`'s \st and \ul are built on a character-by-character
% scan that XeLaTeX's font machinery breaks: the document fails with an
% undefined control sequence rather than degrading. ulem does the same two jobs
% and is Unicode-engine safe. `normalem` stops it hijacking \emph, which the
% transcriptions use for Grothendieck's own emphasis.
\usepackage[normalem]{ulem}
\usepackage{hyperref}
\hypersetup{colorlinks=true,linkcolor=black,urlcolor=black,pdfborder={0 0 0}}

% --- Métadonnées du lot -----------------------------------------------------
% Reprises telles quelles par le script de rendu, qui les affiche en tête de
% la vue de lecture. Elles fixent ce que le fichier prétend couvrir : sans
% elles, une transcription flotte sans point d'attache dans un fonds de seize
% mille feuillets.

\newcommand{\folder}[1]{\gdef\grothFolder{#1}}
\newcommand{\batch}[1]{\gdef\grothBatch{#1}}
\newcommand{\leaves}[2]{\gdef\grothFirst{#1}\gdef\grothLast{#2}}
\newcommand{\foldertitle}[1]{\gdef\grothFoldertitle{#1}}
\gdef\grothFolder{?}\gdef\grothBatch{?}\gdef\grothFirst{?}\gdef\grothLast{?}\gdef\grothFoldertitle{}

% --- Le résumé --------------------------------------------------------------
%
% Il ouvre la lecture modernisée, sous le titre « Résumé », et il est composé
% en petit. Ce n'est pas un ornement : c'est le seul passage du document qui
% ne soit pas des mathématiques, et un lecteur doit pouvoir le reconnaître
% avant de l'avoir lu — puis le sauter s'il connaît déjà le sujet, ou s'y
% arrêter s'il ne le connaît pas. Le filet qui le suit marque la reprise du
% corps.

\newenvironment{resume}
  {\small\setlength{\parskip}{0.45em}}
  {\par\medskip\hrule\medskip}

% --- Mots-clés ---------------------------------------------------------------
%
% En anglais, en clôture du résumé de la lecture modernisée : le vocabulaire
% moderne sous lequel chercher ce que les feuillets construisent. Le manifeste
% du site (`npm run manifest`) les extrait pour étiqueter la cote entière —
% cette ligne est la source unique des tags, il n'y a pas de fichier à côté.
\newcommand{\keywords}[1]{\par\smallskip\noindent
  {\footnotesize\textsc{Keywords} --- \textit{#1}}\par}

% --- L'appareil critique ----------------------------------------------------

% Le numéro de feuillet, en marge. C'est l'ancre qui permet de régler un
% désaccord sur une formule en regardant *un* feuillet plutôt qu'une chemise
% de six cents. Le site s'en sert aussi pour tourner les pages du fac-similé
% quand on fait défiler la transcription.
\newcommand{\leaf}[1]{%
  \marginnote{\footnotesize\color{gray!70}\textsf{#1}}%
  \label{leaf:#1}\ignorespaces}

% Illisible : on ne devine pas. Un mot inventé qui se lit comme les autres est
% le pire résultat possible.
\newcommand{\ill}{\textcolor{red!60!black}{[\,\ldots\,]}}

% Lecture douteuse : le mot est donné, et signalé comme incertain.
\newcommand{\uncertain}[1]{\uline{#1}}

% Ajout de l'éditeur — une lettre manquante, un mot sous-entendu.
\newcommand{\add}[1]{\textcolor{blue!50!black}{[#1]}}

% Biffé par Grothendieck lui-même. Ce qu'il a rayé fait partie du document :
% une rature dit souvent où la pensée a changé de direction.
\newcommand{\struck}[1]{\sout{#1}}

% Note du transcripteur, sur l'état matériel du feuillet.
\newcommand{\note}[1]{\par\smallskip{\small\color{gray!60!black}\textit{[#1]}}\par\smallskip}

% Note marginale de Grothendieck, distincte du corps du texte.
\newcommand{\marginal}[1]{\par\smallskip\noindent\hspace{1em}%
  {\small\color{gray!60!black}#1}\par\smallskip}

% --- Titre ------------------------------------------------------------------

\AtBeginDocument{%
  \begin{center}
    {\large\bfseries Fonds Alexandre Grothendieck --- cote n\textsuperscript{o}~\grothFolder}\\[2pt]
    {\itshape\grothFoldertitle}\\[4pt]
    {\small feuillets \grothFirst--\grothLast\ (lot \grothBatch)}\\[2pt]
    {\footnotesize\color{gray!60!black}Transcription établie d'après le fac-similé numérisé
     par l'Université de Montpellier.\\
     Les lectures incertaines sont soulignées, les illisibles notées [\,\ldots\,],
     les ajouts entre crochets.}
  \end{center}
  \vspace{1em}}

\endinput


\folder{148}
\batch{4}
\pages{61}{71}
\dating{1983}
\watermark{Édition de démonstration}
\foldertitle{Teichmüller (1983) : notes manuscrites (1983)}

\begin{document}

\section{Le dévissage de $S\mathcal{T}D^{\pm}_{\mathrm{top}}$ (suite)}

\note{ce titre est de l'édition. Les pages 61 à 64 continuent une suite
paginée par lui, commencée dans le lot 3 : p.~61 = sa p.~18. Les opérations
numérotées 1° à 7° sont celles de sa « Récapitulation des opérations
élémentaires », sa p.~17 (p.~60)}

\page{61}
\note{p.~18 de l'auteur}
Il faudrait expliciter un \emph{petit} théorème de dévissage préliminaire,
qui dise qu'à équivalence près (définie à isom.\ unique près)
$S_{f}\mathcal{T}D^{\pm}_{\mathrm{top}}$,\note{l'indice sous le $S$, lu $f$,
est douteux ; il revient p.~61 sous
$S_{f}\mathcal{T}^{\pm}_{\mathrm{top}}$} \add{avec les foncteurs
$\underline{\Gamma}$ dans (Gr.pond.\ top) et $f$ dans $[\pm 1]$, et}\note{ajout
interlinéaire de sa main, au-dessus de la ligne} avec les opérations
élémentaires 1° : 7°\marginal{à l'exception de 4° et 5°}\note{note
marginale oblique, soulignée, reliée à « 1° : 7° »} (avec leurs variantes
composées simultanées) se reconstitue par la connaissance des groupoïdes et
foncteurs suivants

\note{un premier diagramme, biffé en entier, commençait par
$S_1\mathcal{T}^{\pm}_{g,\nu} \to$ \struck{\ill{}} $[\pm 1]$, avec une flèche
verticale marquée (*) vers $\mathcal{T}^{\pm}_{g,\nu}$ et une flèche
oblique, toutes deux barrées ; il est remplacé par la ligne encadrée
suivante}

\begin{tikzcd}
S_1\mathcal{T}^{\pm}_{g,\nu} \arrow[r] & \mathcal{T}^{\pm}_{g,\nu} \arrow[r] & (\mathrm{Ens}(\nu) \times [\pm 1])
\end{tikzcd}

(ou simplement $S\mathcal{T}^{\pm}_{g,\nu}$)\note{sous
$S_1\mathcal{T}^{\pm}_{g,\nu}$, de sa main ; la ligne du diagramme est
encadrée par deux traits horizontaux}

\marginal{\uncertain{Ça ne suffit pas} tout à fait, pour avoir toutes les
opérations. Notamment les opérations (4°) (qui demande
\struck{\ill{}} l'\uncertain{introd.}\ des structures ; \ill{} \ill{}
$\mathcal{T}_{g,\nu}$) et (5°)}\note{note écrite à 45° dans la marge
gauche ; « 4° » et « 5° » y sont cerclés}

$S_1\mathcal{T}^{\pm}_{g,\nu}$ désigne le sous-groupoïde de
$S_{f}\mathcal{T}^{\pm}_{\mathrm{top}}$ formé des $(X, R)$ tels que $X$ soit
connexe de genre $g$, $\operatorname{card} \pi_0\,\partial X = \nu$, et
$\operatorname{card} R \cap C_a = 1$ pour tout $a \in \pi_0\,\partial X$ --- i.e.\
le graphe \add{bipondéré} des \struck{$\mathcal{T}_{g,\nu}$}
$(X, \widetilde{\Sigma}, R)$ est isom.\ à\note{« bipondéré » en interligne ;
sous $\widetilde{\Sigma}$ il écrit « $\partial X$ » avec un signe
d'égalité vertical}

\note{ici un graphe dessiné : un sommet marqué $g$ d'où part un faisceau de
$\nu$ arêtes (« $\nu$ » au bout du faisceau)}

toutes ces arêtes (lisières) sont de poids $\mathbf{1}$

et $\mathcal{T}^{\pm}_{g,\nu}$ est défini de façon correspondante, via les
$(X, \emptyset, \emptyset)$ avec $X$ du type $g, \nu$.

\page{62}
\note{p.~19 de l'auteur}
La question essentielle \ill{} \ill{} \uncertain{comment} reconstituer, pour
un graphe \struck{\ill{}} pondéré : \struck{à repères donné} \add{repéré}
(disons) $\underline{\Gamma}$, le groupoïde
$\mathcal{T}^{\pm}_{\mathrm{top}}(\underline{\Gamma})$, \struck{plus}
\add{avec} sa dépendance fonctorielle en $\underline{\Gamma}$,
\struck{\ill{} foncteurs des} \add{et} l'explicitation des sept opérations, à
l'exception de 4°) et 5°).
\marginal{NB on \uncertain{laisse} ici tomber les « D », « S » \ill{} dans
\ill{} des $\underline{\Gamma}$}\note{note oblique dans la marge gauche,
« NB » souligné ; lecture très incertaine}

Pour l'opération 4°), elle sera obtenue automatiquement par la description
explicite des $\mathcal{T}^{\pm}_{\mathrm{top}}(\underline{\Gamma})$ (qui sera
faite en même temps et de la même façon que pour les
$\mathcal{T}^{\pm}_{g,\nu}$ et $S\mathcal{T}^{\pm}_{g,\nu}$
connus),\note{devant le premier $\mathcal{T}^{\pm}_{g,\nu}$ un petit signe
en exposant, non lu} et les bouclages de trous \struck{s'expliciteront}
\add{pourrait s'expliciter} à l'aide \add{des cas} \struck{des opérations}
particuliers de ces opérations pour les \struck{catégories, ou}
$S\mathcal{T}_{g,\nu}$ et $\mathcal{T}_{g,\nu}$, qu'il faudrait jeter dans le
sac \struck{des} \struck{\ill{}} \add{structures} « élémentaires » :
diminution

\begin{tikzcd}
(**)\quad S\mathcal{T}^{\pm}_{g,\nu,\nu'} \arrow[r] & S\mathcal{T}_{g,\nu\,?\,\nu'}
\end{tikzcd}

\note{les deux $S$ sont traversés de traits verticaux, et une note à droite,
séparée par un trait, demande : « avec ou sans $S$ ? ». Le signe entre $\nu$
et $\nu'$ dans l'indice du but, un I barré, n'est pas lu (on attendrait $+$ ou
$-$)}

\marginal{avec ou sans $S$ ?}

$\mathcal{T}^{\pm}_{g,\nu,\nu'}$ : les $X \in \mathrm{Ob}\,\mathcal{T}^{\pm}_{g,\nu}$,
avec \struck{\ill{}} donnée d'une sous-ens.\ $I' \subset I(X)$ avec
$\operatorname{card} I' = \nu'$.\note{écrit sous une accolade qui porte sur
la source de (**)}

\page{63}
\note{p.~20 de l'auteur}
Mais ces opérations de bouclage sont relativement superficielles, en ce
qu'elles ne permettent pas de relier les $\mathcal{T}_{g,\nu}$, pour $g$
fixé, à des $\mathcal{T}_{g',\nu}$ avec $g' < g$. Pour ceci, il faut
l'opération (5°) de gommage des cercles de découpage (= ex-cercles de
recollement !).\note{« (5°) » est cerclé et ajouté au-dessus de la ligne.
Tel qu'écrit : dans la récapitulation de sa p.~17 (p.~60, lue à basse
résolution) le gommage paraît être le 4°, le bouclage le 5°}

Néanmoins, il y a, \struck{en} \uncertain{d'ailleurs}, une autre structure
importante à expliciter, c'est celle des \struck{\ill{}} groupoïdes
fondamentaux \add{$\Pi_1(X; R)$} des $(X, [\widetilde{\Sigma}], R)$,
\add{pour $R \neq \emptyset$,} points-base dans $R$, \add{en tant que}
groupoïdes dépendant fonctoriellement bien sûr de $R$, et
\struck{la dernière} \ill{} \ill{} \ill{} opérations) pour les
\struck{\ill{}} \add{\ill{}} \ill{} groupoïde par \uncertain{ces} finis,
\add{\struck{\ill{}} \ill{} d'un} \uncertain{revêtements} $X'$ des $X$),
d'un nouvel objet des \struck{\ill{}} groupoïdes
$S\mathcal{T}^{\pm}_{\mathrm{top}}$. Déjà dans le cas de
$S_2\mathcal{T}_{0,3}$ la question se pose et est non triviale --- en fait,
c'est sans doute là une exigence trop ambitieuse, car\note{la fin de la phrase
entre « et » et « d'un nouvel objet » est chargée de ratures et d'ajouts
interlinéaires ; l'ordre donné ici est incertain}
\marginal{$\widetilde{\Sigma}$ ici \ill{} \struck{\ill{}} \ill{}, \ill{}
\ill{} les groupoïdes \ill{} dans les \ill{} \ill{} \ill{} (i.e.\
$S\mathcal{T}^{\pm}_{\mathrm{top}} \subset S\mathcal{T}D^{\pm}_{\mathrm{top}}$)
\ill{}}\note{note oblique dans la marge gauche, en face des lignes sur
$\Pi_1$}

\page{64}
\note{p.~21 de l'auteur, la dernière de la suite}
elle exigerait d'avoir un « modèle » pour $S\mathcal{T}^{\pm}_{\mathrm{top}}$
qui soit assez gros pour contenir (\ill{} \ill{}) \emph{toutes les cordes}
--- ça fait un peu gros ! Mais quand même on doit pouvoir sans trop de mal
déterminer dans le \ill{}, \ill{} le \uncertain{modèle} \ill{} \ill{} pour
$S\mathcal{T}D^{\pm}_{\mathrm{top}}$, des \add{modèles des} $\Pi_1(X, R)$, et
leur dépendance fonctorielle en $(X, R)$\ldots

\section{Découpage de cartes en « pantalons » généralisés}

\note{ce titre est de sa main, au crayon, seul en haut à droite de la page 65,
qui sert de couverture à ce qui suit}

\page{66}
\note{feuille d'esquisses, au crayon et à l'encre ; un bloc de texte, en haut,
est écrit par-dessus des dessins}
Un pantalon \add{(topologique)} subdivisé de façon standard, avec un ens.\
$J$ de trous ($\operatorname{card} J = 3$) correspond à un « triangle
pondéré » i.e.
\[
J \xrightarrow{\ \pi\ } \mathbb{N}
\]
(NB $\pi(j)$ = nb d'arêtes qui relient $j$ et $k$, où $i \neq j \neq k$)
\note{la parenthèse est de lui ; « relient $j$ et $k$ » est lu tel quel}

Carte \add{orientée} \struck{découpée} \uncertain{pantalonniquement} de façon
standard $\Longleftrightarrow$ \struck{\ill{}} \uncertain{une} carte
pantalonnière généralisée \struck{i.e.\ carte cellulaire (pas néc.\
orientée), dont les sommets} \ill{} sont d'ordre 3 (sauf des sommets
\uncertain{topologiques} \ill{} d'ordre 1 éventuellement), et avec une
pondération des \struck{arêtes repérées} \add{\ill{}} (sommet et arête
incidente) de telle façon que l'on ait pour chaque arête (cf.\ ci-contre)
\[
\alpha_s + \beta_s = \alpha_t + \beta_t
\]
\note{sous le membre de gauche, une accolade marquée $\gamma'_s$, sous le
membre de droite une accolade marquée $\gamma'_t$. À côté, une arête $st$
dessinée, avec les poids $\alpha_s, \beta_s, \gamma_s$ en $s$ et $\alpha_t,
\beta_t, \gamma_t$ en $t$}
et $=$ de données d'une \struck{\ill{}} \supplied{anti-}isom.\ des torseurs
$(\gamma = \gamma'_s = \gamma'_t)$ sous $\mathbb{Z}/\gamma$\,: sommes
circulaires de $[1, \alpha_s]$ et $[1, \beta_s]$ $\rightleftarrows$ sommes
circulaires de $[1, \alpha_t]$ et $[1, \beta_t]$.\note{« $\mathbb{Z}/\gamma$ »
douteux (l'indice est surchargé) ; la flèche entre les deux « sommes
circulaires » est rendue par l'édition, il y a sur la page un signe barré. Au
bout de la ligne « et avec une » un ajout, « \ill{} de types $0, 1$ », non
situé}

\[
\begin{array}{lll}
\beta + \gamma = \alpha' & \qquad & \alpha = \dfrac{\beta' + \gamma' - \alpha'}{2} = \dfrac{\alpha' + \beta' + \gamma'}{2} - \alpha' \\[2ex]
\gamma + \alpha = \beta' & & \beta = \dfrac{\gamma' + \alpha' - \beta'}{2} \\[2ex]
\alpha + \beta = \gamma' & & \gamma = \dfrac{\alpha' + \beta' - \gamma'}{2}
\end{array}
\]
\note{la dernière ligne de droite est surchargée (« $-(\beta-\gamma)$ » biffé à
côté) ; au crayon, en dessous, un essai
« $2(\alpha+\beta+\gamma) = \ldots$, $2\alpha + \beta + \gamma = \ldots$,
$\beta + \gamma = \alpha'$ »}

\emph{NB} $\alpha, \beta, \gamma \in \mathbb{N}$ ssi $\alpha' + \beta' +
\gamma'$ pair\note{« $\in \mathbb{N}$ » est lu sur un signe rapide ; la ligne
peut porter une autre condition}

\begin{itemize}
  \item trois sommets, $\alpha + \beta + \gamma$ \struck{bigônes} arêtes
  \item \uncertain{3} sommets d'indice $\alpha$, $\beta$, $\gamma$
  \item $(\alpha - 1) + (\beta - 1) + (\gamma - 1)$ $= \alpha + \beta + \gamma - 2$ faces-bigônes, 2 faces-triangles
  \item $\alpha + \beta + \gamma - 1$ faces
\end{itemize}
\note{deux colonnes séparées par un trait vertical : à gauche les deux
premières et la troisième ligne, à droite la deuxième et la quatrième. Devant
la troisième, un « $\alpha+\beta+\gamma$ » biffé}
\note{à côté d'une sphère dessinée au crayon avec trois sommets $0$, $1$,
$\infty$ reliés par des faisceaux d'arêtes. Le « $= \alpha+\beta+\gamma-2$ »
coiffe par une accolade la somme $(\alpha-1)+(\beta-1)+(\gamma-1)$ ; le dernier
chiffre, lu 2, est douteux (la somme vaut $\alpha+\beta+\gamma-3$). Le mot
biffé devant « arêtes » est lu « bigônes » sous réserve}

\note{Dessins de la page. À gauche, à l'encre rouge et noire, la sphère à trois
trous $0$, $1$, $\infty$ subdivisée : sur le bord du trou $0$ les points
$a^0_1, \ldots, a^0_4$ et $b^0_1, \ldots, b^0_5$, sur le trou $1$ les points
$a^1_1, a^1_2, a^1_3$ et $b^1_1, \ldots, b^1_4$, sur le cercle extérieur
$\infty$ les points $a^\infty_1, a^\infty_2, a^\infty_3$ et $b^\infty_1,
\ldots, b^\infty_5$, reliés deux à deux par des arcs. En dessous, des cercles
orientés, dont l'un marqué « $[0, \alpha][1, \beta)$ » (lecture douteuse), et
deux triangles de sommets $0$, $1$, $\infty$ : le premier avec $\alpha = 3$
(points $a_1, a_2, a_3$ sur le côté $1\infty$), $\beta = 5$ ($b_1, \ldots,
b_5$ sur $0\infty$), $\gamma = 4$ ($c_1, \ldots, c_4$ sur $01$) ; le second
avec la suite « $c_1 c_2 c_3 c_4 \cdot b_5 b_4 b_3 b_2 b_1$ ». Au centre, un
petit triangle subdivisé ($b_1, \ldots, b_5$, $a_1, a_2, a_3$, $c_1, \ldots,
c_4$) ; à droite, une carte trivalente et les régions $A_{0\infty}$, $A_{01}$,
$A_{10}$, $A_{1\infty}$, $A_{\infty 0}$, $A_{\infty 1}$ au crayon}

\page{67}
\note{au dos d'un listing d'ordinateur (imprimé à l'envers, daté « 02 JUN 82 »),
seulement des esquisses de sa main au crayon : une boucle à un sommet, une
arête à deux sommets entourée de courbes, un cercle portant six sommets reliés
par des cordes, et un disque à deux trous avec des segments radiaux --- les
figures de la page 66}

\page{68}
\note{un seul dessin, sans texte : une carte au crayon sur la sphère, dont les
sommets sont entourés de petits cercles marqués $0$, $1$, $\infty$, et une
famille de courbes fermées emboîtées, à l'encre rouge, qui contournent ces
sommets --- un découpage en « pantalons »}

\page{69}
\note{la page est tournée d'un quart de tour ; à gauche, un petit dessin
(un secteur de disque subdivisé), au centre un décalque au crayon du dessin de
la page 68. Trois longs traits obliques au crayon traversent le diagramme
suivant ; ils appartiennent au dessin et ne paraissent pas le biffer}

\begin{tikzcd}[column sep=large]
{[\pm 1]} & {[\pm 1]} \\
\mathcal{T}^{\pm}_{\mathrm{top}} \arrow[u] \arrow[r, hook, "\mathrm{imm}"] & S\mathcal{T}S^{\pm}_{\mathrm{top}} \arrow[u] \\
\mathcal{T}_{\mathrm{top}} \arrow[u, "{\text{fid.\ ess.\ surj.}}"] \arrow[r, hook, "\mathrm{imm}"] & \mathcal{T}S_{\mathrm{top}} \arrow[u, "{\text{fid.\ ess.\ surj.}}"']
\end{tikzcd}

\note{le premier $S$ de $S\mathcal{T}S^{\pm}_{\mathrm{top}}$ est lu sous
réserve}

On a
\[
\pi_0(\mathcal{T}_{\mathrm{top}}) \xrightarrow{\ \sim\ } \pi_0(\mathcal{T}^{\pm}_{\mathrm{top}}) \simeq
\]
\note{la ligne s'arrête là}

\section{Groupoïdes de Teichmüller}

\note{titre de sa main, en tête de la page 70}

\page{70}
1) Groupoïde topologique $S\mathcal{T}_{\mathrm{top}}$\note{encadré} (\emph{groupoïde
de Teichmüller spécial orienté topologique})

Objets : paires $(X, S)$
\begin{itemize}
  \item $X$ surface à bord compacte \emph{orientée}
  \item $S \subset \partial X$ partie finie du bord
\end{itemize}

Morphismes $(X, S) \to (X', S')$ : classes d'isotopie d'isomorphismes.

1°) $S = \emptyset$ on écrit $\mathcal{T}_{\mathrm{top}}$ (\emph{groupoïde de
T.\ \struck{\ill{}} orienté topologique}) sous-groupoïde
\uncertain{ouvert}.

\struck{Vari} \emph{Variantes} 1°) $S\mathcal{T}^{\pm}_{\mathrm{top}}$\note{encadré}
mêmes objets, mais plus de flèches : classes d'isotopie d'homéom.\
$(X, S) \xrightarrow{\ \sim\ } (X', S')$, ne respectant pas l'orientation
nécessairement.

On a
\[
S\mathcal{T}^{\pm}_{\mathrm{top}} \xrightarrow{\ \varepsilon\ } [\pm 1]
\]
\note{au-dessus de la flèche, une lettre biffée et surchargée ; le but est
écrit sur une rature}
(orientation), groupoïde \struck{\ill{}} défini par le groupe $\pm 1$, noté
$[\pm 1]$ ; le \struck{noyau} (= catégorie fibre) est
$S\mathcal{T}_{\mathrm{top}}$.

2°) $(X, \Sigma, S)$, avec
\begin{itemize}
  \item $X$ compacte à bord orientée
  \item $\Sigma$ sous-variété \struck{fermée} \add{compacte} de dim 1 de $X
  \setminus \partial X$
  \item $S \subset \Sigma \cup \partial X$ partie finie
\end{itemize}

morphismes : classes d'isotopie d'homéom.\ respectant l'orientation.
Notation $\mathcal{T}D_{\mathrm{top}}$ ((D) initiale de « divisage »
\struck{« \ill{} »}) : « \emph{groupoïde de Teichmüller orienté divisé
topologique} ».\note{le $D$ de la notation est écrit sur une rature}

\emph{NB} $S\mathcal{T}_{\mathrm{top}} \subset S\mathcal{T}D_{\mathrm{top}}$
(\add{imm.\ ouverte} \struck{\ill{} pleine}) $=$ sous-groupoïde des objets
tels que $\Sigma = \emptyset$.

3°) On définit de même $\mathcal{T}S^{\pm}_{\mathrm{top}}$ (\emph{groupoïde de
Teichmüller biorienté divisé topologique})

\note{la page est écrite par-dessus un décalque, au crayon, du dessin de la
page 68}

\page{71}
\note{page de dessin, qui revient au découpage en pantalons des pages 66-69}
découpage pantalonnesque d'une carte (triangulée)\note{légende de sa main, en
haut à droite}

\note{à gauche, une triangulation au crayon (neuf sommets entourés de petits
cercles) et, à l'encre rouge, neuf courbes fermées emboîtées qui la découpent
en pantalons ; à droite, un graphe trivalent pondéré, un hexagone avec des
arêtes pendantes et une boucle, les poids $9$, $19$, $3$, $7$, $2$, $1$, $5$,
$4$, $15$, $3\,6$, $4$, $2$, $17$, $3$, $3$ inscrits aux sommets et aux arêtes,
et trois poids cerclés (1), (2), (3)}

\end{document}
